%!TEX program = lualatex
% Arquivo: Template_Estat_Texto.tex

\documentclass[a4paper,12pt]{article}

% =========================================================
% Idioma e tipografia (LuaLaTeX)
% =========================================================
\usepackage[portuguese]{babel}
\usepackage{fontspec}
\setmainfont{Latin Modern Roman}

% =========================================================
% Layout
% =========================================================
\usepackage[top=3cm,bottom=2.5cm,left=3cm,right=2.5cm]{geometry}
\usepackage{setspace}
\onehalfspacing
\usepackage{indentfirst}
\setlength{\parindent}{1.25cm}

% =========================================================
% Matemática (align*, etc.)
% =========================================================
\usepackage{amsmath}
\usepackage{amssymb}

% =========================================================
% Tabelas, listas, figuras
% =========================================================
\usepackage{booktabs}
\usepackage{array}
\usepackage{enumitem}
\setlist[itemize]{noitemsep,topsep=3pt}
\setlist[enumerate]{noitemsep,topsep=3pt}

\usepackage{graphicx}
\usepackage{float}
\usepackage{pdflscape}
\graphicspath{{./}{./imagens/}{./figuras/}}

% =========================================================
% Links
% =========================================================
\usepackage{hyperref}
\usepackage{xurl}
\hypersetup{
  colorlinks=true,
  linkcolor=black,
  urlcolor=blue,
  citecolor=black,
  pdfauthor={},
  pdftitle={}
}

% =========================================================
% Texto de preenchimento (lorem)
% =========================================================
\usepackage{lipsum}

% =========================================================
% Dados da capa
% =========================================================
\newcommand{\Instituicao}{Universidade do Estado do Rio de Janeiro / UERJ}
\newcommand{\Curso}{Labofis / Oceanografia}
\newcommand{\Professor}{Professor: Alexandre}
\newcommand{\Aluno}{Aluno: José Mauro Xavier Elsas}
\newcommand{\Cidade}{Rio de Janeiro}
\newcommand{\Ano}{2026}

\title{\textbf{Procedimento de Manutenção//
Bóia AXIS}}
\author{\Aluno}
\date{\Cidade --- \Ano}

\begin{document}

% =========================================================
% Capa (página de rosto)
% =========================================================
\begin{titlepage}
\centering
{\large \Instituicao \par}
\vspace{0.5cm}
{\large \Curso \par}
\vspace{1.0cm}

{\LARGE \textbf{Procedimento de Manutenção\\ 
Boia AXYS}\par}
\vspace{0.6cm}
\vspace{0.6cm}
{\Large Poseidon 1750 WKB. \par}

\begin{figure}[H]
  \centering
  \includegraphics[width=.3\linewidth]{./uerj.jpg}
\end{figure}

\vfill



\begin{flushright}
\Aluno\\
\Professor
\end{flushright}

\vspace{1.0cm}
{\Cidade --- \Ano \par}
\end{titlepage}

% =========================================================
% Sumário
% =========================================================
\tableofcontents
\newpage

% =========================================================
% Corpo do texto
% =========================================================

\section{Escopo}

Esta ficha técnica preliminar foi montada a partir dos desenhos e manuais disponíveis para a boia AXYS WatchKeeper / Poseidon 1750 WKB.

Foram analisadas duas configurações principais:

\begin{itemize}
  \item Modelo 4-BAT: banco com 4 baterias Rolls AGM 12 V / 116 Ah.
  \item Modelo 6-BAT: banco com 6 baterias Rolls AGM 12 V / 116 Ah.
\end{itemize}

O objetivo é consolidar os dados relevantes para planejamento técnico, manuseio, desenho do procedimento de recuperação (se necessário), estimativa preliminar de massa para içamento e identificação das pendências necessárias para fechamento da ficha técnica final.

\subsection{Documentos Consultados}

\begin{itemize}
  \item AXYS 424101 — WatchKeeper Buoy Operator, Installation and User Manual.
  \item AXYS 424000 — WKB, 1.75 m Hull Buoy, Standard Options, General Specifications.
  \item AXYS 9517 — 1.75 m Hull Buoy, Basic without Options, Assembly Drawing.
  \item AXYS 9100 — 1.75 m Hull Battery Assembly.
\end{itemize}

\subsection{Resumo executivo}

As direções de corrente, fornecidas em graus azimutais, foram - O modelo 4-BAT possui peso seco indicado de 650 kg sem bridle e 750 kg com bridle.

\begin{itemize}
  \item O modelo 6-BAT não aparece diretamente como “6-BAT WKB sem bridle” na tabela de massas, mas pode ser estimado de forma consistente como 710 kg sem bridle.
  \item O peso seco estimado do modelo 6-BAT com bridle é 810 kg.
  \item O SWL/WLL da cinta de içamento não foi identificado nos documentos disponíveis.
  \item Para procedimento operacional, o SWL/WLL deve ser obtido no certificado da sling, da manilha ou do conjunto de rigging efetivamente usado na operação.
   Para desenho preliminar de recuperação, recomenda-se usar 810 kg como massa seca de referência do conjunto 6-BAT com bridle, sem confundir esse valor com a capacidade admissível da sling.

Os valores extremos foram tratados por meio de um filtro baseado no percentil 99 da magnitude para cada célula, removendo apenas outliers associados a ruído instrumental. A comparação entre resultados com e sem filtragem indicou impacto inferior a 5\%, garantindo a robustez da análise.

\end{itemize}

\section{Identificação Geral do equipamento}

Equipamento:
AXYS WatchKeeper Buoy / Poseidon 1750 WKB.

Casco:
Poseidon 1750, diâmetro nominal de 1,75 m.

Configuração 4-BAT:
Banco com 4 baterias Rolls AGM 12 V / 116 Ah.

Configuração 6-BAT:
Banco com 6 baterias Rolls AGM 12 V / 116 Ah.

Aplicação:
Boia oceanográfica / meteoceanográfica com sensores, telemetria e alimentação solar.


\subsection{Dimensões e Características Mecânicas Mecânicas comuns}

Diâmetro do casco:
1,75 m.

Diâmetro da superestrutura:
0,4 m no topo do top cone; 0,7 m no flange da mid-section.

Altura total:
4,7 m, da base ao topo das antenas.

Altura com instrumentação removida:
3,65 m.

Material do casco:
Polietileno amarelo, espessura 9,5 mm, densidade 920 kg/m³.

Lastro:
Concreto, 330 kg.

Enchimento do casco:
Espuma de poliuretano de célula fechada.

Fixadores:
Aço inox SS316.

Estrutura dos painéis solares:
Alumínio AL 6061 com pintura eletrostática.

Estrutura interna:
Alumínio AL 6061.

Balde das baterias:
Alumínio AL 6061 + aço inox.

\subsection{Pesos principais}

Configuração 4-BAT:
\begin{itemize}
  \item Peso seco sem bridle: 650 kg.
  \item Peso seco com bridle: 750 kg. 
  \item Observação: valor indicado nos documentos.
\end{itemize}

Configuração 6-BAT:
\begin{itemize}
\item Peso seco sem bridle: 710 kg.
\item Peso seco com bridle: 810 kg.
\item Observação: valor inferido pelo acréscimo de 2 baterias de 30 kg cada.
\end{itemize}

\section{Cálculo  de peso adotado para o modelo 6-BAT}

Cálculo adotado para o modelo 6-BAT.

\begin{table}[H]
\centering
\begin{tabular}{cc}
\toprule
Descrição  & Valor (kg)
\midrule
Peso Seco sem bridle & 650\\
Peso Seco com bridle & 750\\
\bottomrule
\end{tabular}
\caption{Resumo da Configuração}
\end{table}

\section{Fundamentação teórica}
\lipsum[3]

\subsection{Um exemplo com \texttt{align*}}
\begin{align*}
f(x) &= -0.5x^2 + 2x - 3 \\
g(x) &= -0.5x - 0.5
\end{align*}

\lipsum[4]

\section{Metodologia}
\lipsum[5]

\section{Resultados}
\lipsum[6]

% Exemplo de tabela
\begin{table}[H]
\centering
\caption{Exemplo de tabela com \texttt{booktabs}.}
\label{tab:exemplo}
\setlength{\tabcolsep}{12pt}
\renewcommand{\arraystretch}{1.2}
\begin{tabular}{lc}
\toprule
\textbf{Intervalo} & \textbf{Frequência} \\
\midrule
$85 \le x < 100$  & 5 \\
$100 \le x < 115$ & 5 \\
\bottomrule
\end{tabular}
\end{table}

\begin{table}[H]
\centering
\begin{tabular}{cccc}
\toprule
Célula & Direção (°) & Velocidade (mm/s) & Observação \\
\midrule
001 & 182.8 & 106.7 & Superfície \\
002 & 188.1 & 154.2 &  \\
003 & 190.5 & 194.7 &  \\
004 & 192.6 & 204.2 & Máximo \\
005 & 194.5 & 200.6 &  \\
006 & 196.0 & 186.7 &  \\
007 & 198.4 & 173.7 &  \\
008 & 201.0 & 160.9 &  \\
009 & 203.0 & 144.4 &  \\
010 & 205.2 & 130.0 &  \\
\bottomrule
\end{tabular}
\caption{Vetor médio por célula (após filtragem p99).}
\end{table}

\end{document}
